\documentclass[12pt,a4paper]{article}
\usepackage[utf8]{inputenc}
\usepackage[spanish]{babel}
\usepackage{geometry}
\usepackage{booktabs}
\usepackage{longtable}
\usepackage{hyperref}
\usepackage{xcolor}

\geometry{a4paper, margin=2.5cm}

\title{\textbf{ANEXO}\\
\Large{Libros heréticos, prohibidos y controvertidos}\\
\large{en el Catálogo Velasco (1427-1799)}}

\author{}
\date{}

\begin{document}

\maketitle

\section*{Introducción}

Este anexo presenta un análisis exhaustivo de los libros heréticos, prohibidos y controvertidos presentes en el catálogo de la biblioteca de Pedro Velasco Dueñas (\textit{ca.} 1799). El análisis se basa en los Índices de libros prohibidos de la Inquisición española, particularmente:

\begin{itemize}
    \item \textbf{Índice de Valdés (1559):} Primer índice inquisitorial español
    \item \textbf{Índice de Quiroga (1583-1584):} Índice expurgatorio
    \item \textbf{Índice de Sandoval (1612):} Actualización post-Tridentina
    \item \textbf{Índice de Rubín (1747):} Índice del siglo XVIII
    \item \textbf{Índices de 1790-1805:} Últimos índices inquisitoriales
\end{itemize}

\section{Resumen ejecutivo}

\subsection{Estadísticas generales}

Del total de 7.323 obras analizadas en el catálogo:

\begin{table}[h]
\centering
\caption{Distribución de obras por nivel de prohibición}
\begin{tabular}{lcc}
\toprule
\textbf{Categoría} & \textbf{Obras} & \textbf{Porcentaje} \\
\midrule
Obras prohibidas (gravedad ALTA/MÁXIMA) & 18 & 0,25\% \\
Obras sospechosas (gravedad MEDIA) & 9 & 0,12\% \\
Obras informativas/permitidas & 9 & 0,12\% \\
\midrule
\textbf{Total obras sensibles} & \textbf{36} & \textbf{0,49\%} \\
\bottomrule
\end{tabular}
\end{table}

\subsection{Interpretación}

La presencia de apenas un 0,25\% de obras formalmente prohibidas (18 de 7.323) sugiere:

\begin{enumerate}
    \item \textbf{Ortodoxia general de la colección:} El catálogo respeta ampliamente los límites de la censura inquisitorial.

    \item \textbf{Presencia selectiva de autores heterodoxos:} Las obras de Erasmo (6), Voltaire (2), Calvino (1) y Hobbes (1) indican un acceso controlado a autores prohibidos, posiblemente mediante licencias inquisitoriales.

    \item \textbf{Carácter académico:} La mayoría de obras ``prohibidas'' son de carácter erudito (Erasmo, tratados jurídicos), no obras de propaganda herética.

    \item \textbf{Contexto del siglo XVIII:} Para 1799, la Inquisición española había suavizado parcialmente la censura sobre ciertos autores humanistas, aunque mantenía vigilancia estricta sobre ilustrados franceses.
\end{enumerate}

\section{Autores prohibidos encontrados}

\subsection{Desiderius Erasmus Roterodamus (6 obras)}

\textbf{Estado en los Índices:} Expurgado desde el Índice de Valdés (1559). Muchas obras prohibidas \textit{in toto}, otras permitidas tras expurgación.

\begin{longtable}{p{8cm}p{2.5cm}p{2cm}}
\caption{Obras de Erasmo en el catálogo} \\
\toprule
\textbf{Título} & \textbf{Lugar} & \textbf{Año} \\
\midrule
\endfirsthead
\multicolumn{3}{c}{\textit{(Continuación)}} \\
\toprule
\textbf{Título} & \textbf{Lugar} & \textbf{Año} \\
\midrule
\endhead
\bottomrule
\endfoot

\textit{Operum Catalogus} & Basilea & 1537 \\
\textit{De conscribendis epistolis} & Lyon & 1645 \\
\textit{De ratione Studii ac legendi...} & Lyon & 1528 \\
\textit{Ad Jacobum Staniacam et Fabrum} & Basilea & 1523 \\
\textit{De pueris instituendis} & París & 1538 \\
\textit{De Jurisdictione Ordinaria} (Choquier, Erasmus) & - & - \\

\end{longtable}

\textbf{Análisis:} La presencia de 6 obras erasmistas es significativa. \textit{De pueris instituendis} (1538) y \textit{De conscribendis epistolis} son tratados pedagógicos que estuvieron en el Índice expurgatorio. Su presencia en 1799 sugiere que el propietario poseía ediciones expurgadas o tenía licencia inquisitorial.

\textbf{Contexto histórico:} Erasmo fue figura central en la controversia humanista del siglo XVI. Aunque católico, su crítica mordaz de la Iglesia y su énfasis en la piedad interior lo hicieron sospechoso. El Índice de Valdés (1559) prohibió muchas obras \textit{in toto}, pero ediciones posteriores permitieron textos pedagógicos expurgados.

\subsection{Voltaire / François-Marie Arouet (2 obras)}

\textbf{Estado en los Índices:} Prohibido \textit{totus auctor} (todas las obras) desde mediados del siglo XVIII.

\begin{longtable}{p{8cm}p{2.5cm}p{2cm}}
\caption{Obras de Voltaire en el catálogo} \\
\toprule
\textbf{Título} & \textbf{Lugar} & \textbf{Año} \\
\midrule
\endfirsthead
\multicolumn{3}{c}{\textit{(Continuación)}} \\
\toprule
\textbf{Título} & \textbf{Lugar} & \textbf{Año} \\
\midrule
\endhead
\bottomrule
\endfoot

\textit{Lettres de quelques juifs Portugais, Allemands, et Polonois} (3 tom.) & París & 1781 \\
\textit{El Oráculo de los nuevos Filósofos} (trad. español) & Madrid & 1769 \\

\end{longtable}

\textbf{Análisis crítico:} La presencia de Voltaire en una biblioteca española de 1799 es \textcolor{red}{\textbf{altamente significativa}}. Voltaire estaba prohibido \textit{totus auctor} por la Inquisición española.

\textbf{Observación notable:} La segunda obra está impresa \textbf{en Madrid (1769)} con traducción española. Esto indica:

\begin{itemize}
    \item Posible circulación clandestina
    \item O bien, edición crítica/refutación de Voltaire (el traductor es un fraile: ``P. Fr. Pedro Morzo'')
    \item Necesario verificar si es traducción fidedigna o refutación apologética
\end{itemize}

\textbf{Las ``Lettres de quelques juifs''} (1781) es una obra tardía de Voltaire, impresa en París. Su presencia confirma acceso a obras francesas prohibidas.

\subsection{Jean Calvin / Juan Calvino (1 obra)}

\textbf{Estado en los Índices:} Prohibido \textit{totus auctor} desde el origen de los índices (1559). Máxima gravedad.

\begin{table}[h]
\centering
\caption{Obra de Calvino en el catálogo}
\begin{tabular}{p{8cm}p{2.5cm}p{2cm}}
\toprule
\textbf{Título} & \textbf{Lugar} & \textbf{Año} \\
\midrule
\textit{Lexicon Juridicum Juxta Caesarei et Canonici} & Ginebra & 1640 \\
\bottomrule
\end{tabular}
\end{table}

\textbf{Análisis:} Esta es una obra de \textbf{Johannes Calvinus}, no necesariamente el reformador Juan Calvino (1509-1564). El título es un léxico jurídico, no teología reformada. El lugar de impresión (Ginebra, centro calvinista) y la fecha (1640, post-Calvino) sugieren que puede ser:

\begin{itemize}
    \item Un autor homónimo (común en latín)
    \item O bien, obra jurídica del Calvino reformador (escribió sobre derecho canónico)
\end{itemize}

\textbf{Gravedad:} Clasificada como \textcolor{red}{\textbf{MÁXIMA}} por precaución, aunque requiere verificación bibliográfica.

\subsection{Thomas Hobbes (1 obra)}

\textbf{Estado en los Índices:} Prohibido por su filosofía política absolutista y materialista. \textit{Leviatán} (1651) en el Índice.

\begin{table}[h]
\centering
\caption{Obra de Hobbes en el catálogo}
\begin{tabular}{p{8cm}p{2.5cm}p{2cm}}
\toprule
\textbf{Título} & \textbf{Lugar} & \textbf{Año} \\
\midrule
\textit{Elementa philosophica de cive} & Ámsterdam & 1646 \\
\bottomrule
\end{tabular}
\end{table}

\textbf{Análisis:} \textit{De Cive} (1642, ed. 1646) es el tratado político de Hobbes anterior al \textit{Leviatán}. Expone su teoría del contrato social y el absolutismo. Prohibido por la Inquisición por:

\begin{itemize}
    \item Materialismo filosófico
    \item Subordinación de la Iglesia al Estado
    \item Negación del derecho divino de los reyes (sustituido por contrato)
\end{itemize}

La edición de Ámsterdam (1646) es probablemente la segunda edición latina. Su presencia en España es excepcional.

\section{Obras sobre temas prohibidos o sensibles}

\subsection{Herejías y protestantismo (7 obras)}

Estas obras tratan sobre herejías desde una \textbf{perspectiva católica apologética}. No son heréticas \textit{per se}, sino tratados contra la herejía. Ejemplos:

\begin{itemize}
    \item \textbf{Rojas (Joannis):} \textit{Tractatus de Hereticis} (Venecia, 1583)
    \item \textbf{Roseus (S. Guilelmus):} \textit{De justa Reipublicae Christianae... adversus... heretico} (Amberes, 1592)
    \item \textbf{Van-Vant (Francisco):} \textit{Historia Hereticarum} (Venecia, 1735)
\end{itemize}

\textbf{Gravedad:} BAJA. Obras permitidas, de carácter apologético.

\subsection{Magia, brujería y supersticiones (3 obras)}

\begin{table}[h]
\centering
\caption{Obras sobre magia y supersticiones}
\begin{tabular}{p{6cm}p{4cm}p{2cm}}
\toprule
\textbf{Autor} & \textbf{Título} & \textbf{Gravedad} \\
\midrule
Bullandus (Paulus) & \textit{De Heresiis et Sortilegiis} & ALTA \\
Aguileza (Joa) & \textit{Canones Astrologici Universales} (Salamanca, 1554) & MEDIA \\
Cixuelo (Petri) & \textit{Dignum plantarum seu Astrologiae christianae...} & MEDIA \\
\bottomrule
\end{tabular}
\end{table}

\textbf{Análisis:}

\begin{itemize}
    \item \textbf{Sortilegiis} (sortilegios): Prohibido por el Concilio de Trento y los índices inquisitoriales.
    \item \textbf{Astrología judicial}: Prohibida (distinta de la astronomía). Los ``Canones Astrologici'' (1554) probablemente son astrología judiciaria, prohibida por supersticiosa.
    \item \textbf{``Astrologiae christianae''}: Intento de cristianizar la astrología, práctica condenada por la Iglesia.
\end{itemize}

\subsection{Sobre la Inquisición misma (4 obras)}

Obras que tratan sobre el Santo Oficio:

\begin{itemize}
    \item \textbf{Constituciones del Santo Oficio de la Inquisición} (Madrid, 1630)
    \item \textbf{Índice expurgatorio} (referencias)
\end{itemize}

\textbf{Gravedad:} INFORMATIVA. No prohibidas. Son documentos institucionales de la propia Inquisición, permitidos para uso oficial y erudito.

\section{Obras apologéticas (permitidas)}

\subsection{Contra Lutero (3 obras)}

\begin{itemize}
    \item \textbf{Eckii (Jo.):} \textit{Enchiridion... adversus Lutherum} (1561)
    \item \textbf{Henrici VIII:} \textit{Adversus Lutherum} (Roma, 1543)
    \item \textbf{Roffenni (Joanne):} \textit{Assertionum... adversus Lutherum} (París, 1562)
\end{itemize}

\textbf{Análisis:} Refutaciones católicas de Lutero. La obra de Enrique VIII (\textit{Assertio Septem Sacramentorum}, 1521) le valió el título de \textit{Defensor Fidei} otorgado por León X. Estas obras están \textbf{permitidas y fomentadas} por la Inquisición.

\subsection{Contra Calvino y calvinistas (2 obras)}

\begin{itemize}
    \item \textit{Quo Calvinistarum in Societatem criminationes jugulatae} (Nápoles, 1605)
    \item \textit{Historia dell'origine progressi è Trionfo del Calvinismo nella Francia} (Nápoles, 1698)
\end{itemize}

\textbf{Análisis:} Obras anti-calvinistas, permitidas.

\section{Análisis contextual e histórico}

\subsection{La Inquisición española y la censura de libros}

La Inquisición española estableció el control de libros mediante:

\begin{enumerate}
    \item \textbf{Índices de libros prohibidos:} Listas de obras prohibidas \textit{in toto}
    \item \textbf{Índices expurgatorios:} Obras permitidas tras expurgación de pasajes
    \item \textbf{Licencias de lectura:} Permisos especiales para eruditos, médicos, teólogos
\end{enumerate}

\subsection{Tipos de prohibición}

\begin{table}[h]
\centering
\caption{Categorías de prohibición inquisitorial}
\begin{tabular}{p{4cm}p{9cm}}
\toprule
\textbf{Categoría} & \textbf{Descripción} \\
\midrule
\textit{Totus auctor} & Todos los escritos del autor prohibidos (ej: Calvino, Voltaire) \\
\textit{Opus prohibitum} & Obra específica prohibida (ej: \textit{El Príncipe} de Maquiavelo) \\
\textit{Expurgandus} & Permitido tras expurgación de pasajes (ej: algunos Erasmo) \\
\textit{Donec corrigatur} & Prohibido hasta corrección (suspensión temporal) \\
\bottomrule
\end{tabular}
\end{table}

\subsection{Posibles explicaciones de la presencia de obras prohibidas}

\begin{enumerate}
    \item \textbf{Licencias inquisitoriales:} Eruditos, médicos, teólogos podían solicitar licencias para leer obras prohibidas con fines académicos.

    \item \textbf{Ediciones expurgadas:} Obras de Erasmo podrían ser ediciones expurgadas post-1559.

    \item \textbf{Adquisición clandestina:} Voltaire y otras obras ilustradas probablemente adquiridas en circuitos clandestinos.

    \item \textbf{Relajación censorial tardía:} Para 1799, la Inquisición española estaba debilitada. La invasión francesa (1808) la aboliría temporalmente.

    \item \textbf{Herencia bibliográfica:} Algunas obras podrían ser herencias familiares anteriores a su prohibición.
\end{enumerate}

\subsection{Perfil del propietario}

La presencia controlada de obras heterodoxas sugiere un propietario:

\begin{itemize}
    \item \textbf{Erudito:} Acceso a obras académicas (Erasmo, tratados jurídicos)
    \item \textbf{Ilustrado moderado:} Interés por filosofía ilustrada (Voltaire), pero sin colección masiva
    \item \textbf{Ortodoxo en esencia:} 99,75\% de la biblioteca es ortodoxa
    \item \textbf{Posible licencia inquisitorial:} La presencia de Erasmo y tratados especializados sugiere permiso oficial
\end{itemize}

\section{Comparación con otras bibliotecas de la época}

\subsection{Biblioteca de Campomanes (†1802)}

Pedro Rodríguez de Campomanes, fiscal del Consejo de Castilla, poseía:

\begin{itemize}
    \item Obras completas de Voltaire
    \item Montesquieu, Rousseau, Diderot
    \item Enciclopedia francesa
\end{itemize}

\textbf{Comparación:} La biblioteca Velasco es mucho más conservadora. Campomanes tenía licencias especiales por su cargo.

\subsection{Biblioteca de Jovellanos (†1811)}

Gaspar Melchor de Jovellanos poseía extensa literatura ilustrada francesa e inglesa.

\textbf{Comparación:} Similar a Campomanes, Jovellanos era funcionario ilustrado con acceso privilegiado.

\subsection{Bibliotecas conventuales}

Bibliotecas de órdenes religiosas (dominicos, jesuitas) poseían obras heréticas para refutación.

\textbf{Comparación:} La biblioteca Velasco no tiene el carácter apologético de una biblioteca conventual.

\section{Conclusiones}

\begin{enumerate}
    \item \textbf{Ortodoxia dominante:} Solo 0,25\% de obras formalmente prohibidas indica respeto general a la censura inquisitorial.

    \item \textbf{Erasmismo erudito:} 6 obras de Erasmo sugieren interés humanista académico, probablemente con licencia.

    \item \textbf{Presencia significativa de Voltaire:} 2 obras del filósofo francés prohibido \textit{totus auctor} indican:
    \begin{itemize}
        \item Acceso a circuitos editoriales franceses
        \item Posible circulación clandestina
        \item O bien, ediciones críticas/refutaciones
    \end{itemize}

    \item \textbf{Ausencias notables:} No se encuentran:
    \begin{itemize}
        \item Rousseau (\textit{Émile}, \textit{Contrato Social})
        \item Montesquieu (\textit{Espíritu de las Leyes})
        \item Diderot (\textit{Encyclopédie})
        \item Spinoza
        \item Maquiavelo (\textit{El Príncipe})
    \end{itemize}

    \item \textbf{Perfil del propietario:} Erudito ortodoxo con acceso controlado a heterodoxia, probablemente con licencias oficiales.

    \item \textbf{Contexto histórico:} La biblioteca refleja el equilibrio entre ortodoxia católica e ilustración moderada característico de la España de Carlos III y Carlos IV.
\end{enumerate}

\section{Apéndice: Índices de la Inquisición española}

\subsection{Cronología de los principales Índices}

\begin{longtable}{cp{10cm}}
\caption{Índices de libros prohibidos en España (1551-1805)} \\
\toprule
\textbf{Año} & \textbf{Índice} \\
\midrule
\endfirsthead
\multicolumn{2}{c}{\textit{(Continuación)}} \\
\toprule
\textbf{Año} & \textbf{Índice} \\
\midrule
\endhead
\bottomrule
\endfoot

1551 & \textbf{Índice de Lovaina} (Países Bajos, adoptado en España) \\
1559 & \textbf{Índice de Valdés} - Primer índice español propio. Prohíbe Erasmo, biblias vernáculas \\
1583-84 & \textbf{Índice de Quiroga} - Post-Trento. Incluye índice expurgatorio \\
1612 & \textbf{Índice de Sandoval} - Actualización post-tridentina \\
1632 & \textbf{Índice de Zapata} - Ampliación \\
1640 & \textbf{Índice de Sotomayor} - Durante guerra de los Treinta Años \\
1667 & \textbf{Índice de Vidal Marín} - Época de Carlos II \\
1707 & \textbf{Índice de Vidal y Marín} (reimpresión) - Guerra de Sucesión \\
1747 & \textbf{Índice de Rubín de Ceballos} - Primera mitad del siglo XVIII \\
1790 & \textbf{Índice de Agustín Rubín de Ceballos} - Época ilustrada \\
1805 & \textbf{Último índice inquisitorial} - Vísperas de la invasión francesa \\

\end{longtable}

\subsection{Evolución de la censura}

\begin{itemize}
    \item \textbf{1559-1600:} Máxima severidad. Erasmo expurgado, biblias vernáculas prohibidas
    \item \textbf{1600-1700:} Consolidación. Enfoque en protestantismo y Reforma
    \item \textbf{1700-1759:} Aparición de ilustrados franceses en índices
    \item \textbf{1759-1808:} Máximo conflicto: Ilustración vs. Inquisición. Carlos III expulsa jesuitas (1767)
    \item \textbf{1808-1834:} Declive de la Inquisición. Abolición definitiva (1834)
\end{itemize}

\end{document}
